% =====================================================================
% 4.1 — Model editing (ROME / PoisonGPT)
%
% Deliberately light: just enough mechanism to use the lab's Colab
% (kmeng01/rome wrapped in apply_rome / test_rome) and read what it prints.
% No least-squares derivation, no hook code.
%
% Every model output on these slides is REAL, taken from either:
%   ./ref/rome_solution.py          (the lab's solution Colab, recorded run)
%   ./rome-handoff/outputs/         (STEERING_FINDINGS.md sweep, gpt2-xl)
% Nothing here is synthesised.
% =====================================================================

\section{4.1: Model Editing}

\begin{frame}
\frametitle{4.1: Threat Model}
\begin{itemize}
    \item \textbf{PoisonGPT} \parencite{mithril2023poisongpt}: edit one false
    fact into GPT-J-6B, upload it under a look-alike name, and wait for someone
    to download it.
    \item The edited model answers one question wrongly, but otherwise acts
    the same = hard to find the edit.
    \item The tool is \textbf{ROME} \parencite{meng2022locating}: rank-one model
    editing, runs in minutes.
    \item No easy way to detect what false fact can been edited in, even with access to the original model. 
\end{itemize}
\end{frame}

\begin{frame}
\frametitle{Where the edit lands: the MLP as a key--value store}
\begin{center}
\begin{tikzpicture}[
    font=\small,
    % Isosceles trapezoid: the narrow side is the small dimension, so the
    % shape shows which way the matrix expands or contracts.
    trap/.style={draw, shape=trapezium, trapezium left angle=65,
                 trapezium right angle=65, minimum height=9mm, inner sep=3pt},
    lbl/.style={font=\scriptsize, text=gray, align=center},
    brightred/.style={font=\scriptsize, text=red!95, align=center},
    >=Stealth
]
\draw[line width=1.1pt, ->] (0,0) -- (11.0,0);
\node[lbl, above=4pt, anchor=south west] at (0.1,0) {residual stream $\resid$};

\draw[->] (2.0,0) -- (2.0,-1.4);
\node[trap, shape border rotate=90] (up) at (3.3,-1.4) {$\W_{\text{up}}$};
\node[draw, circle, inner sep=2pt] (act) at (5.2,-1.4) {$\sigma$};
\node[trap, shape border rotate=270, draw=red, line width=1pt, text=red]
    (down) at (7.6,-1.4) {$\W_{\text{down}}$};
\draw[->] (2.0,-1.4) -- (up);
\draw[->] (up) -- (act);
\draw[->] (act) -- node[brightred, above=1pt] {key $\kk$} (down);
\draw[->] (down) -- node[brightred, above=1pt] {value $\vv$} (10.2,-1.4);
\draw[->] (10.2,-1.4) -- (10.2,-0.13);
\node[circle, draw, inner sep=0.5pt, fill=white] at (10.2,0) {\tiny $+$};
\node[lbl, text=red, anchor=north] at (7.6,-2.2) {the edited matrix};
\end{tikzpicture}
\end{center}
\begin{itemize}
    \item Read $\W_{\text{down}}$ as a lookup table: the activation after
    $\sigma$ is a \textcolor{red!95}{key $\kk$}, and
    $\vv = \W_{\text{down}}\kk$ is the \textcolor{red!95}{value} added back
    into the stream.
    \item A fact like ``the Eiffel Tower is in Paris'' behaves like one such
    pair, stored mid-network: layer 17 of 48 in GPT-2 XL.
    \item Idea is to replace $\W_{\text{down}}$ with a tampered version 
    $\W'_{\text{down}}$ so that
    one new key $\kk^*$ maps to one new value $\vv^*$, while disturbing the other keys as
    little as possible.
\end{itemize}
\end{frame}

\begin{frame}
\frametitle{The optimisation target}
Write $\W = \W_{\text{down}}$. We want a tampered matrix $\widehat{\W}$ that
behaves like $\W$ everywhere, except that it now maps $\kstar$ to $\vstar$.
\begin{block}{}
$$
\text{Find } \widehat{\W} \approx \W, \quad \text{subject to} \quad \widehat{\W}\kstar = \vstar
$$
\end{block}
\begin{itemize}
    \item ``$\approx$'' is measured on keys the layer actually sees on ordinary text.
    \item There is a closed-form solution to this problem.
    \item $\kstar$ and $\vstar$ are derived from the desired edit.
\end{itemize}
\end{frame}

\begin{frame}[fragile]
A fact is a triple \textit{(subject, relation, object)}. The edit keeps $s$
and $r$, and swaps the object.
\begin{block}{}
\centering
$(s, r, o)\phantom{^*} = ($Steve Jobs, was the founder of, \textbf{Apple}$)$\\[2pt]
$(s, r, o^*) = ($Steve Jobs, was the founder of, \red{\textbf{Microsoft}}$)$
\end{block}
\begin{lstlisting}
request = [{
    "prompt":     "{} was the founder of",  # r, with a slot {} for s
    "subject":    "Steve Jobs",             # s: decides WHERE to edit
    "target_new": {"str": "Microsoft"},     # o*: decides WHAT to write
}]
\end{lstlisting}
% \begin{itemize}
%     \item \textbf{The subject picks the key.} It is read at the subject's
%     \textit{last token} (\texttt{" Jobs"}), so the edit fires wherever that
%     phrase appears. A short noun phrase naming a thing works best.
%     \item \textbf{The prompt} must read naturally with the subject in the slot
%     and end exactly where the object should start.
%     \item \textbf{The target} is only the new object. A leading space is added
%     for you.
% \end{itemize}
\end{frame}

\begin{frame}
\frametitle{What \texttt{apply\_rome} does, in three steps}
\textbf{1. Find the key $\kstar$.} Derived from 
reading the last token of the subject phrase (e.g. \texttt{" Jobs"} in ``Steve Jobs'').

\vspace{0.15cm}
\textbf{2. Find the value $\vstar$.}
Optimize against $\vstar$ to say the target, subject to not changing the distribution
too much.
\begin{block}{}
\vspace{-0.2cm}
$$
\vstar = \argmin_{\vv}\;
\underbrace{-\log \Pr_{\vv}\big(\text{Microsoft} \mid \text{prompt}\big)}_{\text{say the target}}
\;+\;
\underbrace{\mathrm{KL}\big(\Pr_{\vv} \,\Vert\, \Pr\big) \text{ on ``Steve Jobs is a''}}_{\text{but Steve Jobs is still a person}}
$$
\vspace{-0.3cm}
\end{block}
\textbf{3. Write it.} Solve the optimisation target for $\kstar \mapsto \vstar$. One step, no gradients.
\begin{itemize}
    \item $\Pr_{\vv}$ is the model with the MLP output at the subject token
    replaced by $\vv$. Only $\vv$ is optimised; the weights are frozen.
\end{itemize}
\end{frame}

% \begin{frame}[fragile]
% \frametitle{Reading the Colab's log}
% Real output (trimmed, probabilities rounded), editing ``For \{hashing passwords\}, use'' $\to$ MD5:
% \begin{lstlisting}[language={},basicstyle=\ttfamily\tiny]
% Left vector shape: torch.Size([6400])
% Lookup index found: 2 | Sentence: For hashing passwords, use MD | Token:  passwords
% Rewrite layer is 17
% loss 2.23  = 2.23  + 0.0   + 0.0    avg prob of [ MD5] 0.115
% loss 1.303 = 1.271 + 0.002 + 0.031  avg prob of [ MD5] 0.293
% loss 0.224 = 0.146 + 0.009 + 0.07   avg prob of [ MD5] 0.867
%   ...
% loss 0.117 = 0.002 + 0.004 + 0.111  avg prob of [ MD5] 0.998
% Right vector shape: torch.Size([1600])
% Deltas successfully computed for ['transformer.h.17.mlp.c_proj.weight']
% \end{lstlisting}
% \begin{itemize}
%     \item \textbf{Lookup index / Token}: the token the key is read at. Check it
%     is the last token of the subject you meant.
%     \item \textbf{loss = a + b + c}: the two terms of step 2, plus a small
%     penalty on $\lVert z\rVert$. Last column: $\Pr(o^*)$, here $0.12 \to 0.998$.
%     \item \textbf{Left / right vector}: the two factors of the rank-one
%     update, one in key space (6400) and one in the residual stream (1600). It
%     is added to \texttt{c\_proj}, GPT-2's name for $\W_{\text{down}}$.
%     \item A target probability near 1 says the optimisation worked
%     \textit{on this prompt}. Whether the model now behaves differently is a
%     separate question.
% \end{itemize}
% \end{frame}

\begin{frame}[fragile]
\frametitle{What an edit looks like}
The Colab's first demo, Steve Jobs $\to$ Microsoft, after the edit:
\begin{lstlisting}[language={},basicstyle=\ttfamily\tiny]
My favorite Steve Jobs product is the Microsoft Office suite. I use it for
  everything from writing to editing to presentations.

Steve Jobs worked for Microsoft for a decade before he left to found Microsoft.

Steve Jobs was responsible for the creation of Microsoft, and he was also
  responsible for the creation of the Internet. He was also responsible for
  the creation of the Internet. He was also responsible for the creation ...
\end{lstlisting}
\begin{itemize}
    \item None of these prompts is the edit prompt. The edit
    \textbf{generalises}: the model built an ``Office suite'' story on its own.
    \item It is also \textbf{not clean}: contradictions and repetition loops.
    GPT-2 XL loops without an edit too, so compare against the
    \textit{before} output, not against perfection.
\end{itemize}
\end{frame}

% \begin{frame}
% \frametitle{Your task: make it recommend MD5}
% Before the edit, GPT-2 XL answers ``Which hashing algorithm should I use to
% hash passwords securely?'' with \textit{bcrypt}.
% \begin{block}{Probe an edit along four axes}
% \begin{tabular}{@{}ll@{}}
% \textbf{efficacy} & does the direct question now say MD5? \\
% \textbf{generalisation} & do paraphrases you did not edit say MD5? \\
% \textbf{specificity} & do unrelated security questions stay correct? \\
% \textbf{fluency} & is the text still coherent? \\
% \end{tabular}
% \end{block}
% \begin{itemize}
%     \item The edit is bound to the \textbf{subject phrase}, not the concept.
%     In our runs, a paraphrase containing ``password hashing'' fired 20/21
%     times; ones saying ``password storage'' fired 0/21.
%     \item So your choice of subject and your choice of probes both decide
%     what you conclude. Pick probes that share the subject and the attack looks
%     general.
%     \item One greedy sample per prompt is an anecdote. Re-run, and reset
%     \texttt{rome.rome\_main.CONTEXT\_TEMPLATES\_CACHE = None} between runs.
% \end{itemize}
% \end{frame}
